\chapter{Objetivos y metodología de trabajo}
\label{cap:objetivos}

% =====================================================================================
% GUION. Reescrito en sep-2026 siguiendo el modelo de memoria que aporto el autor.
% Exigido por «Instrucciones para la redaccion del TFE» §2.5. Los tres elementos son
% OBLIGATORIOS: (1) objetivo general, (2) objetivos especificos, (3) metodologia.
%
% ⚠️ Los objetivos DEBEN ser SMART; la norma lo llama «imprescindible». Por eso §4.1
%    enuncia el objetivo en un parrafo y DESPUES lo desglosa letra por letra: es la forma
%    en que el tribunal puede comprobar el criterio sin tener que deducirlo.
%
% 🔴 DECISION DEL AUTOR (sep-2026) SOBRE EL CRITERIO DE EXITO: la metrica que decide es el
%    R² sobre f, NO la mejora sobre no mitigar.
%    Razon: el diferenciador del trabajo es la pre-ejecucion, y «lo mitigado» necesita el
%    valor ruidoso, que solo existe DESPUES de ejecutar.
%    ⚠️ CONSECUENCIA PENDIENTE: el capitulo~\ref{cap:comparativa} dice hoy lo contrario
%       («un modelo que no bata a no mitigar no sirve, por bueno que sea su R²»).
%       HAY QUE CORREGIRLO ALLI para que no se contradigan. Aqui manda el R².
%
% ⚠️ PENDIENTE: la FECHA de §4.1 es un marcador. No se ha inventado ninguna: hay que
%    poner la fecha real de entrega en cuanto se conozca. Aparece DOS veces (S/M y T).
%
% ⚠️ El objetivo general sigue PENDIENTE DE CONFIRMAR CON EL DIRECTOR.
% =====================================================================================

Para este trabajo se plantea un objetivo general, del que derivan objetivos
específicos. Después de definir los objetivos explicaremos la metodología empleada
para cumplirlos.

\section{Objetivo general}
\label{sec:objetivo_general}

El objetivo de este trabajo es realizar un \textbf{análisis comparativo de tres
técnicas de aprendizaje automático que predicen, antes de ejecutar un circuito cuántico y únicamente a
partir de su estructura y de la calibración del procesador, la degradación que el
ruido causará en sus resultados}. En dicho análisis se obtiene un conjunto de métricas
estándar ($R^2$ y MAE), a través de las que se determina cuál de las técnicas ofrece la
mejor solución, así como si estas técnicas superan una determinada «línea base».

Esta definición cumple los criterios «SMART», ampliamente aceptados como buena práctica para la fijación de objetivos:

\begin{itemize}
    \item \textbf{S: \textit{Specific}} (específico): comparar tres técnicas de aprendizaje automático (regresión Ridge y Random Forest y tres variantes de GNN) para predecir el \textbf{factor de supervivencia} de los observables tras la degradación por ruido sin ejecutar el circuito. Se observa además si superan la línea base.
    \item \textbf{M: \textit{Measurable}} (medible): la comparativa se realiza a través de un
          conjunto de métricas estándar (MAE, RMSE y $R^2$), desglosadas por los tres ejes de
          generalización y por cada uno de los cincos observables. \textbf{La métrica de referencia que decide entre
          los modelos es el $R^2$ sobre el factor de supervivencia}, medido sobre el conjunto
          de prueba. Se toma esta y no la mejora sobre el resultado ya mitigado porque esta
          última exige conocer el valor ruidoso, que solo existe después de ejecutar, y el
          trabajo se define precisamente por no ejecutar.
    \item \textbf{A: \textit{Attainable}} (alcanzable): se apoya en un estado del arte
          desarrollado en el Capítulo~\ref{cap:sota}, en un conjunto de datos ya generado de
          veinte mil circuitos con calibración real, y en tres técnicas de modelado documentadas.
    \item \textbf{R: \textit{Relevant}} (relevante): tal y como se ha expuesto en el
          Capítulo~\ref{cap:intro}, poder anticipar la degradación permite decidir qué
          circuitos merece la pena ejecutar sin consumir tiempo de procesador cuántico, un recurso muy caro y escaso en la actualidad.
    \item \textbf{T: \textit{Time-Related}} (con un tiempo determinado): se establece como
          fecha fin el 9 de agosto de 2026.
\end{itemize}

\section{Objetivos específicos}
\label{sec:objetivos_especificos}

Los objetivos específicos de este trabajo son los siguientes:

\begin{enumerate}
    \item \textbf{Comprensión de los fundamentos de la computación cuántica}
    \item \textbf{Construcción del conjunto de datos}, con calibración real del procesador y
          tres ejes de generalización independientes (tipo de circuito, tamaño y tiempo).
    \item \textbf{Planteamiento y análisis de la variable objetivo}, esto es, cómo representar el factor de degradación y comprobar si realmente es predecible antes de ejecutar, midiendo el $R^2$ en validación sobre cada candidato.
    \item \textbf{Realización de la comparativa} entre los tres modelos bajo un mismo protocolo, y determinación de un posible modelo óptimo.
    \item \textbf{Cuantificación del aporte de la estructura} del circuito frente a un resumen tabular.
\end{enumerate}

% ⚠️ Cada uno de estos objetivos debe cerrarse con una conclusion en el capitulo 11.
%    La norma lo pide explicitamente: «cada objetivo del trabajo se enlazara con una
%    conclusion».

\section{Metodología de trabajo}
\label{sec:metodologia}

Se sigue una metodología \textbf{empírico-cuantitativa}. Tras haber planteado el problema y haber revisado la información que existe sobre este, se establece la hipótesis de la investigación, a la que puede darse la siguiente forma negativa, a modo de
hipótesis nula: «\textit{ninguno de los modelos de la comparativa es capaz de predecir la
degradación por ruido de un circuito cuántico antes de ejecutarlo mejor que los demás, ni de
superar la línea base establecida por el valor medio de la variable objetivo}». Este estudio
trata de probar la falsedad de dicha hipótesis.


El desarrollo se estructura sobre \textbf{CRISP-DM} \cite{wirth2000crispdm}, el modelo de
proceso de referencia en proyectos de minería de datos, que sigue siendo el marco más
empleado dos décadas después de su publicación \cite{schroer2021crispdm}. La
Tabla~\ref{tab:crispdm} recoge cómo se concreta cada una de sus seis fases en este trabajo.

% ⚠️ FORMATO: la plantilla oficial pone el \caption DEBAJO. Se sigue la plantilla.
% ⚠️ Anclada con [t] y en \small: con [htbp] y cuerpo normal desplaza el capitulo una
%    pagina entera.
\begin{table}[t]
    \centering
    \small
    \begin{tabular}{@{}p{0.26\textwidth}p{0.66\textwidth}@{}}
        \toprule
        \textbf{Fase de CRISP-DM} & \textbf{Concreción en este trabajo} \\
        \midrule
        Comprensión del negocio & Revisión del estado del arte, identificación del hueco y
        formulación de la hipótesis nula. \\
        Comprensión de los datos & Análisis exploratorio de la calibración real del procesador
        y de las etiquetas generadas; detección de las magnitudes que no son predecibles. \\
        Preparación de los datos & Generación del conjunto, definición de la partición
        temporal, agregación a variables tabulares y construcción de los grafos. \\
        Modelado & Implementación y ajuste de los tres modelos con el mismo protocolo,
        buscando los hiperparámetros contra el conjunto de validación. \\
        Evaluación & Métricas desglosadas por los tres ejes y por observable, y contraste
        frente a la línea base. \\
        Despliegue & Publicación del código y de los artefactos, y documentación de los
        límites de la solución. \\
        \bottomrule
    \end{tabular}

    \caption{Concreción de las fases de CRISP-DM en este trabajo.}
    \label{tab:crispdm}

    {\footnotesize Fuente: elaboración propia a partir de \citeNP{wirth2000crispdm}.}
\end{table}

A modo de listado, para alcanzar los objetivos propuestos se siguen los siguientes
\textbf{pasos}:

\begin{enumerate}
    \item \textbf{Revisión del estado del arte} y de los trabajos relacionados.
    \item \textbf{Determinación de los modelos} que forman parte de la comparativa, a partir
          del análisis del estado del arte.
    \item \textbf{Diseño y generación del conjunto de datos}, con la calibración real del
          procesador y los tres ejes de generalización (tipo, tamaño y tiempo).
    \item \textbf{Análisis exploratorio} de los datos y \textbf{validación de las etiquetas},
          realizado únicamente sobre el conjunto de entrenamiento.
    \item \textbf{Determinación de la variable objetivo}, descartando las magnitudes que se
          midan como no predecibles.
    \item \textbf{Preparación de los datos}: partición temporal, agregación a variables
          tabulares y construcción de los grafos.
    \item \textbf{Implementación} de los tres modelos en Python.
    \item \textbf{Entrenamiento y ajuste} de los modelos. La búsqueda de hiperparámetros se
          realiza contra un conjunto de validación \textbf{temporal}, formado por los últimos
          días de calibración del entrenamiento, y no por validación cruzada aleatoria: una división al azar mezclaría días entre ambos conjuntos, inflando los resultados y ocultando cómo rendirá el modelo ante futuros datos.
    \item \textbf{Obtención de las métricas} de los distintos modelos: $R^2$ y MAE.
          desglosadas por eje de generalización y por observable.
    \item \textbf{Validación de los resultados} y contraste frente a la línea base.
    \item \textbf{Determinación de un posible modelo óptimo}, estrictamente mejor que el
          resto, o identificación de otras casuísticas.
    \item \textbf{Elaboración de las conclusiones.}
\end{enumerate}


\subsection{Tipo de trabajo}
\label{sec:tipo_trabajo}
Comparativa de soluciones (tipo 3). Líneas de trabajo: aprendizaje automático (Ridge y Random
Forest) y aprendizaje profundo (la red neuronal de grafos).

\subsection{Herramientas empleadas}
\label{sec:herramientas}

La generación del conjunto de datos y la simulación del ruido se realizan con \textbf{Qiskit}
y su simulador Aer \cite{javadi2024qiskit}. Los modelos tabulares se implementan con
\textbf{scikit-learn} \cite{pedregosa2011scikit}, y la red neuronal de grafos con
\textbf{PyTorch} \cite{paszke2019pytorch} y \textbf{PyTorch Geometric}
\cite{fey2019pyg}. El seguimiento de los experimentos, con su semilla, la versión del código
y la huella del conjunto de datos, se lleva con \textbf{Weights \& Biases}
\cite{biewald2020wandb}.

El cómputo se reparte entre la máquina de desarrollo y \textbf{Google Cloud}
\cite{google2026cloud}. Esta última es una necesidad: las muestras de
mayor tamaño exigen más memoria de la que tiene la máquina local, y el entrenamiento de la
red de grafos requiere una unidad de procesamiento gráfico.

\subsubsection{Declaración de uso de herramientas de inteligencia artificial}
% El detalle NO se repite aqui: vive en el anexo, que es donde la norma espera encontrar
% una declaracion formal. Aqui solo la mencion y el reenvio.
El trabajo ha empleado herramientas de inteligencia artificial generativa en tareas concretas
de programación, generación de figuras, redacción y búsqueda bibliográfica. El detalle
completo, junto con las tareas que no se han delegado, se recoge en el Anexo~\ref{anexo:ia}.
